\chapter{ADVANTAGES AND IMPLEMENTATION CHALLENGES}
\label{ch:advantages_challenges}

\section{Key Advantages of the VibeGuard System}
The engineering architecture of VibeGuard delivers several distinct advantages over conventional commercial and cloud-centric vibration monitoring platforms:

\begin{enumerate}
    \item \textbf{Complete Edge Autonomy and Cybersecurity Resilience:} VibeGuard executes all signal acquisition, statistical filtering, DC gravity cancellation, Vector RMS computation, and spectral peak extraction entirely within the local ESP32 microcontroller. Operating with zero dependence on Wi-Fi, Ethernet, or cloud servers, the system is fundamentally immune to remote network outages, packet dropouts, cloud subscription fees, and cyber attack vectors targeting industrial IoT gateways.
    \item \textbf{Transparent and Explainable Decision Path:} Modern machine learning classifiers often function as opaque black boxes, making it difficult for plant operators to understand why an alert was triggered. VibeGuard relies on first-principles physics and deterministic mathematics: static gravity is canceled via in-place zero-mean subtraction, dynamic vibration severity is quantified via Euclidean Vector RMS, and anomalous states are supported by verifiable FFT spectral peaks at the fundamental shaft rotational frequency ($1X$). Every state transition is traceable, verifiable, and explainable.
    \item \textbf{Ultra-Low Cost and Democratized Accessibility:} Commercial industrial vibration analyzers and piezoelectric accelerometers typically cost between \$1,000 and \$5,000 per monitoring channel. By utilizing commodity MEMS sensors (ADXL345) and an off-the-shelf dual-core microcontroller (7Semi ESP32-DEVKIT-E), VibeGuard achieves complete physical implementation for only \rupee 2,124.00 ($\approx$\$25.00). This drastic cost reduction makes automated condition monitoring accessible to small-scale workshops, educational laboratories, and auxiliary machinery.
    \item \textbf{High Noise Immunity via $K$-of-$M$ Persistence Filtering:} In industrial plant environments, transient mechanical disturbances---such as tool impacts, floor vibrations, or motor startup surges---frequently trigger false alarms on simple thresholding systems. VibeGuard's temporal $K$-of-$M$ persistence filter ($K=4, M=5$) suppresses isolated spikes while maintaining a rapid, deterministic alert detection latency of $1.28$\,seconds for sustained anomalies.
    \item \textbf{Multi-Planar Sensitivity via Triaxial Vector RMS:} Rather than monitoring a single geometric axis, VibeGuard evaluates the composite Euclidean magnitude $a_{\mathrm{vec}} = \sqrt{\mathrm{RMS}_x^2 + \mathrm{RMS}_y^2 + \mathrm{RMS}_z^2}$. This makes the monitoring metric rotationally invariant to mounting orientation while capturing both radial unbalance (in the $X\text{--}Y$ plane) and axial looseness (along the $Z$ axis).
    \item \textbf{Rigorous Data Provenance and Traceability:} The project enforces strict configuration identifiers ($\mathrm{CFG}\text{-}\#\#\#$), linking every recorded telemetry block to its corresponding hardware carrier revision, firmware commit hash, sampling parameters, and mechanical mounting state.
\end{enumerate}

\section{Mechanical and Acoustic Challenges}
Developing a physical, rotating vibration testbed presented several intricate mechanical engineering challenges:

\subsection{Structural Foundation Stiffness and Acrylic Through-Bolting}
Early prototypes utilizing thin wooden boards or hollow 3D-printed baseplates exhibited severe flexural resonance modes below 200\,Hz. The foundation flexed under motor torque, absorbing vibration energy and skewing accelerometer readings. 
\begin{itemize}
    \item \textbf{Resolution:} The baseplate was upgraded to a 10.0\,mm thick laminated optical-grade cast acrylic slab (mass $\approx 265$\,g). Acrylic provides high structural stiffness and acoustic damping.
    \item \textbf{Through-Bolting:} Furthermore, because self-tapping screws in acrylic strip under cyclic shear stresses, all eight mounting holes were precision-drilled completely through the 10\,mm plate and clamped from the bottom using M3 machine screws, wide washers, and nylon-insert locknuts (Nyloc).
\end{itemize}

\subsection{Acoustic Cavity Rattle and Motor Bed Shim (FAB-03)}
Due to the stamped sheet-metal geometry of commercial N20 motor brackets, a 0.8\,mm gap existed between the cylindrical motor casing and the flat acrylic plate. During rotation, the cantilevered motor body vibrated against the bracket wall, creating an acoustic rattling noise at gear-meshing harmonics.
\begin{itemize}
    \item \textbf{Resolution:} A custom 1.0\,mm thick PLA shim pad (\textbf{FAB-03}) was CAD-modeled and 3D-printed to cradle the motor belly. Sandwiched between the motor and the acrylic baseplate, FAB-03 eliminated mechanical play and damped acoustic rattling.
\end{itemize}

\subsection{Centrifugal Separation of 3D-Printed Rotors}
Initial multi-component rotor arm assemblies suffered catastrophic structural failure: under 600\,RPM rotation, dynamic shear forces caused the printed arm to physically delaminate and separate into three disconnected segments.
\begin{itemize}
    \item \textbf{Resolution:} The rotor cam arm (\textbf{FAB-01}) was completely redesigned as a 100\% monolithic, single-piece solid CAD model. Printed with 100\% solid infill and 4 perimeter walls, the monolithic part eliminated layer boundary stress concentrations, successfully withstanding continuous cyclic unbalance forces.
\end{itemize}

\subsection{FDM Thermal Shrinkage Compensation}
Fused Deposition Modeling (FDM) in Polylactic Acid (PLA) inherently shrinks by approximately 1.5\% upon cooling. An as-modeled 3.00\,mm D-shaft bore printed undersized, preventing insertion onto the motor shaft without damaging the motor gearbox.
\begin{itemize}
    \item \textbf{Resolution:} The CAD model incorporated empirical shrinkage compensation, dimensioning the bore to 3.15\,mm diameter and 2.65\,mm flat-to-back chord. This produced a snug, zero-backlash press-fit directly off the print bed without requiring manual filing.
\end{itemize}

\section{Electrical and Sampling Challenges}
\subsection{Inductive Motor Flyback and Electrical Isolation}
Direct-current electric motors generate substantial inductive back-EMF spikes ($L \frac{di}{dt} > 100$\,V) whenever the circuit is switched off. In initial breadboard trials without isolation, motor flyback coupled into the 3.3\,V logic rail, causing microcontroller lockups.
\begin{itemize}
    \item \textbf{Resolution:} Complete electrical domain separation was enforced: the 12\,V motor power loop is powered by an independent 2\,A DC adapter and protected by a fast-acting 1\,A fuse and a reverse-biased \textbf{1N4007 flyback suppression diode} soldered directly across the motor terminals.
\end{itemize}

\subsection{Microcontroller Brownout Resets}
During high-speed SPI transactions and boot sequences, the ESP32's internal current demand can surge transiently to 500\,mA. High internal resistance across thin breadboard jumpers caused the 3.3\,V rail to sag below 2.70\,V, triggering the ESP32 internal Brownout Detector (\texttt{rst:0x10}).
\begin{itemize}
    \item \textbf{Resolution:} A dual-capacitor decoupling network---comprising a \textbf{100\,$\mu$F bulk electrolytic capacitor} and a \textbf{100\,nF low-ESR ceramic capacitor}---was installed directly across the 3.3\,V and GND power rails adjacent to the ESP32, completely stabilizing the voltage rail.
\end{itemize}

\subsection{ESP32 Bootloader Strapping Pin Conflict (GPIO5 Remapping)}
In standard ESP32 reference designs, the VSPI Chip Select pin defaults to GPIO5. However, GPIO5 is a critical boot strapping pin sampled by the ROM bootloader during reset to configure SDIO flash timing. When the ADXL345 pulled GPIO5 LOW during power-up, the MCU entered a bootloop.
\begin{itemize}
    \item \textbf{Resolution:} The active-low Chip Select line was remapped in firmware to \textbf{GPIO21}, an unstrapped general-purpose I/O pin, permanently resolving boot instabilities.
\end{itemize}

\section{Algorithmic and Validation Challenges}
\subsection{Mitigating Temporal Data Leakage}
In vibration condition monitoring, adjacent 320\,ms windows from a single physical run share identical motor temperatures, contact friction, and line voltage. Randomly partitioning windows into training and test sets produces extreme data leakage, yielding artificially inflated accuracy.
\begin{itemize}
    \item \textbf{Resolution:} VibeGuard enforces strict \textbf{session-level isolation}. All windows from a single physical run remain in the same data split. Baseline calibration is performed strictly on initial training sessions, while final validation is conducted on completely independent sessions recorded on separate days or after full mechanical remounting.
\end{itemize}

\subsection{Balancing Real-Time Acquisition with Spectral Computation}
Executing continuous 256-point complex FFTs across three axes at 800\,Hz consumes significant CPU time, risking FIFO overflows and dropped samples during long-term monitoring.
\begin{itemize}
    \item \textbf{Resolution:} The firmware separates continuous time-domain surveillance from spectral analysis. Fast, lightweight Vector RMS is evaluated continuously (consuming only 0.28\,ms per 320\,ms window), while the 256-point Hann-windowed FFT is \textbf{triggered on-demand} upon persistence filter violation. This architecture ensures 100\% acquisition integrity with zero dropped blocks while maintaining total CPU utilization below 0.5\%.
\end{itemize}

\section{Scope Governance and Boundary Discipline}
Throughout the semester execution, strict scope boundaries were maintained. Although machine learning anomaly detection and multi-sensor spatial attribution were explored conceptually, they were intentionally excluded from the semester MVP. Restricting the system to a single rigidly mounted ADXL345, explainable Vector RMS, persistence filtering, and triggered FFT ensured that every reported claim was backed by empirical laboratory evidence rather than speculative simulations.

\section{Summary}
This chapter reviewed the principal advantages and practical engineering challenges of VibeGuard. Key advantages include edge autonomy, cybersecurity resilience, explainable physics-based decision logic, multi-planar sensitivity, and ultra-low implementation cost. Practical challenges spanned mechanical foundation stiffness, 3D printing shrinkage compensation, monolithic rotor redesign, motor flyback suppression, brownout decoupling, and bootloader strapping pin remapping. Finally, session-level validation and on-demand triggered FFT execution resolved the critical issues of temporal data leakage and real-time computational bottlenecks.
